\subsection{Internal Layering of a Service}
\label{subsec:layering}

Every back-end service follows the same four-layer structure, with imports flowing in a single direction:
\begin{equation*}
\texttt{api/} \;\longrightarrow\; \texttt{service/} \;\longrightarrow\; \texttt{repository/} \;\longrightarrow\; \texttt{domain/}
\end{equation*}

\begin{figure}[H]
\centering
\includegraphics[width=0.7\linewidth]{figs/backend_package_diagram.png}
\caption{Package diagram of a typical back-end service.}
\label{fig:backend-pkg}
\end{figure}

The \texttt{api/} layer holds one file per endpoint --- a thin FastAPI router that parses the request, resolves the JWT scope, validates with a Pydantic schema and delegates. The \texttt{service/} layer contains domain logic; it talks only to repositories and never imports SQLAlchemy or FastAPI symbols. The \texttt{repository/} layer is the only one allowed to talk to the database: SQLAlchemy \texttt{text()} queries for analytical reads, SQLModel sessions for CRUD. Dynamic WHERE clauses are built by appending to a \texttt{conditions: list[str]} and a \texttt{params: dict}, joined with \texttt{AND} just before execution --- a pattern that keeps everything parameterised and SQL-injection-safe. The \texttt{domain/} layer is pure data: SQLModel ORM classes and Pydantic schemas.
